\chapter{HCP Coordination and Gravitational Occlusion}
\label{ch:hcp-occlusion}


\section{The Question}

Chapter~\ref{ch:steradian-geometry} derived the inverse-square law from the solid angle subtended by a single sphere. But matter is not a single sphere. At the nuclear scale, matter is an arrangement of nucleon vortices in close-packed configurations. How does the packing geometry of nuclear matter determine the macroscopic gravitational field?

This chapter derives the gravitational occlusion properties of hexagonal close-packed (HCP) and face-centred cubic (FCC) lattices --- the two most common arrangements of nuclear matter at maximum density.


\section{Single Sphere Occlusion}

Consider a sphere of radius $a$ touching a central sphere, also of radius $a$. The angular radius of the touching sphere as seen from the centre of the central sphere:
\begin{equation}
  \sin\alpha = \frac{a}{2a} = \frac{1}{2} \qquad\Rightarrow\qquad \alpha = 30^\circ
\end{equation}

The solid angle subtended by one touching sphere:
\begin{equation}
  \Omega_1 = 2\pi(1 - \cos 30^\circ) = 2\pi(1 - \tfrac{\sqrt{3}}{2}) = 2\pi \times 0.13397 = 0.8418\;\text{sr}
\end{equation}

As a fraction of the full sphere ($4\pi = 12.566$~sr):
\begin{equation}
  f_1 = \frac{0.8418}{12.566} = 6.70\%
\end{equation}

\textbf{Each touching sphere occludes 6.70\% of the full solid angle.}


\section{12-Coordination Shell}

In both HCP and FCC packings, each sphere has exactly 12 nearest neighbours (the \textbf{kissing number} in 3D). The 12 neighbours are arranged in two groups:
\begin{itemize}
  \item \textbf{6 equatorial:} arranged in a hexagonal ring at $\theta = 90^\circ \pm 30^\circ$.
  \item \textbf{3 upper + 3 lower:} arranged in alternating triangular caps.
\end{itemize}

\subsection{Naive Estimate (No Overlap)}

If the 12 occlusion cones did not overlap:
\begin{equation}
  f_{12,\text{naive}} = 12 \times 6.70\% = 80.4\%
\end{equation}

This would leave 19.6\% of the central sphere's solid angle unoccluded.

\subsection{Overlap Correction}

Adjacent spheres in the coordination shell are themselves touching each other. Their occlusion cones overlap at the boundaries. The overlap between two adjacent cones (separated by angular distance $60^\circ$ between centres, each with half-angle $30^\circ$) must be subtracted.

Each sphere touches 5 neighbours within the coordination shell (edges of the triangulation). There are $12 \times 5 / 2 = 30$ overlap pairs. The overlap area per pair (computed from the spherical cap intersection formula) is approximately:
\begin{equation}
  \Delta\Omega_{\text{pair}} \approx 0.0147\;\text{sr}
\end{equation}

Total overlap correction:
\begin{equation}
  \Delta\Omega_{\text{total}} = 30 \times 0.0147 = 0.441\;\text{sr}
\end{equation}

\subsection{Net Occlusion}

\begin{align}
  \Omega_{\text{occluded}} &= 12 \times 0.8418 - 0.441 = 10.10 - 0.44 = 9.66\;\text{sr} \\
  f_{12} &= \frac{9.66}{12.566} = 76.9\%
\end{align}

\textbf{The 12-coordination shell occludes approximately 77\% of the full solid angle.}

The remaining $\sim 23\%$ consists of the ``gaps'' between the 12 touching spheres --- the interstices of the close-packing geometry.


\section{Connection to Gravitational Parameters}

\subsection{Nuclear Packing and Gravitational Strength}

A nucleus with $A$ nucleons packed in HCP geometry creates a composite occlusion pattern. The total gravitational charge of the nucleus depends on:
\begin{enumerate}
  \item The total solid angle occluded by \emph{all} nucleons as seen from infinity (far-field limit).
  \item The packing efficiency of the nuclear lattice.
  \item The nuclear radius: $R_{\text{nuc}} \approx r_0 A^{1/3}$ with $r_0 \approx 1.2$~fm.
\end{enumerate}

In the far-field limit (distance $r \gg R_{\text{nuc}}$), the nucleus appears as a single sphere of radius $R_{\text{nuc}}$, and the occlusion is simply $\Omega = \pi R_{\text{nuc}}^2/r^2$. The internal packing geometry becomes irrelevant at macroscopic distances --- consistent with the observation that gravity depends on total mass (proportional to $A$), not on nuclear structure.

\subsection{Why Packing Matters at Short Range}

At distances comparable to the nuclear radius ($r \sim R_{\text{nuc}}$), the internal packing geometry becomes visible. The non-uniform distribution of nucleon occlusion creates:
\begin{itemize}
  \item \textbf{Angular anisotropies:} The pressure field is not perfectly spherically symmetric at nuclear range.
  \item \textbf{Shell effects:} Completed coordination shells (magic numbers: 2, 8, 20, 28, 50, 82, 126) create particularly stable, symmetric occlusion patterns.
  \item \textbf{The nuclear force:} At inter-nucleon distances ($r \sim 2a$), the occlusion is dominated by the nearest-neighbour geometry, creating the enormously strong short-range ``strong force.''
\end{itemize}


\section{Darkstar Internal Structure}

At maximum density (Darkstar interior, Chapter~\ref{ch:cyclical-universe}), nuclear matter is compressed into perfect HCP close-packing. The packing fraction:
\begin{equation}
  \eta_{\text{HCP}} = \frac{\pi}{3\sqrt{2}} = 0.7405
\end{equation}

This means 74.05\% of space is filled with nucleon vortices and 25.95\% is interstitial spation. The 77\% solid-angle occlusion of the coordination shell is consistent with this: the first coordination shell captures most of the pressure field, leaving $\sim 23\%$ to propagate through the interstices.

A Darkstar is therefore a body where:
\begin{enumerate}
  \item All matter is in HCP close-packing.
  \item The coordination shell occludes $\sim 77\%$ of the pressure field per layer.
  \item Multiple layers rapidly attenuate the remaining flux.
  \item At $\kop = 1$, the cumulative occlusion reaches $100\%$ of the surface escape velocity.
\end{enumerate}


\section{Summary}

\begin{enumerate}
  \item Each touching sphere in an HCP lattice occludes $6.70\%$ of the solid angle ($0.842$~sr).
  \item The 12-coordination shell (kissing number) occludes $\sim 77\%$ after overlap correction.
  \item At macroscopic distances, internal packing is invisible and the nucleus appears as a single sphere.
  \item At nuclear distances, packing geometry determines the ``strong force.''
  \item Darkstars are perfect HCP lattices at maximum density, with cumulative occlusion creating $\kop = 1$.
  \item The close-packing fraction ($74.05\%$) and occlusion fraction ($\sim 77\%$) are geometrically consistent.
\end{enumerate}
